\section{Introduction}
\label{sec:introduction}

Galois connections between lattices offer an appealing mathematical foundation for bidirectional dynamic
program slicing, providing an elegant way of characterising input-output relationships in terms of demand and
sufficiency \cite{perera12a,perera16d,ricciotti17}. The upper adjoint, which characterises forward slicing, is
a \emph{sufficiency} operator, picking out all the outputs which can be computed using only a particular set
of inputs; the lower adjoint, which characterises backward slicing, is a \emph{demand} operator, picking out
the specific inputs needed to compute a particular set of outputs. \citet{perera22} extended this approach to
a setting where slices have complements and the slices of a particular program or value form a Boolean
lattice; they then showed how the De Morgan dual of the upper adjoint can be used to pick out \emph{related
outputs}, namely the other outputs that ``compete'' with a given set of outputs in the sense of demanding some
of the same inputs, and proposed this as a framework for understanding the ``brushing and linking'' feature
commonly used in data visualisation.

However, an implementation of Galois slicing, in particular the backwards analysis, efficient enough to
support this kind of interactive application has proved elusive. To compute the two analyses, implementations
to date have relied on a sequential execution record taking the form of a big-step or small-step trace. The
forward analysis uses the trace to ``replay'' the computation, propagating information about resource
availability from inputs to outputs; the backwards analysis uses the trace to ``rewind'' the computation,
propagating information about demand from outputs to inputs. This overly sequential approach makes the
approach unable to exploit the very independence that makes slicing useful in the first place, and moreover
the backwards analysis needs to make frequent use of a lattice join operation to combine slicing information
computed along different branches of the computation. Joining closure slices in particular becomes expensive
because of the contained environments (which in turn contain closures), and these shortcomings make the
approach prohibitively slow for interactive use.

\todo{To be continued}

\subsection{Contributions}

\secref{preliminaries} establishes some preliminaries. The specific contributions of the paper are then as
follows:
\begin{itemize}
\item \secref{core} defines the core calculus that is the basis of our implementation and gives an operational
semantics that equips every result with a DDG suitable for slicing.
\item \secref{graph} introduces the mathmatical setting of ``Galois connections for a relation'', formalises
DDGs as one possible implementation of such a relation, and defines adjoint forwards and backwards slicing
operators over DDGs.
\item \secref{evaluation} contrasts the graph-based approach with a trace-based implementation, comparing the
overhead of building a trace vs.~a DDG, slicing over both, and the relative ease of implementing the graphical
approach.
\end{itemize}

\noindent \secref{related-work} discusses related work in more detail \secref{conclusion} concludes with a
discussion on limitations and future work.
